\section{Related Work}
\label{sec:related}

\subsection{Graph-based vulnerability detection}
Earlier deep-learning detectors used no graph at all.
VulDeePecker~\citep{Li_2018} sliced programs into token sequences around
library-call sites and classified them with a recurrent network, establishing
both the function-level framing and the practice of evaluating on synthetic
corpora such as the Juliet suite~\citep{juliet2017,6329885}. SySeVR~\citep{9321538}
extended the slicing criterion beyond library calls, and Russell et
al.~\citep{russell2018automatedvulnerabilitydetectionsource} showed that
convolutional models over lexed tokens scale to millions of functions. The move to graphs was
motivated by the argument that token order discards the control and data
relations a vulnerability depends on, a shift documented in surveys of the
area~\citep{harzevili2023surveyautomatedsoftwarevulnerability,shimmi2025aibasedsoftwarevulnerabilitydetection,uddin2025deeplearningaidedsoftware,Guo2024}.

A prevalent methodology represents functions as graphs and classifies them using graph neural networks. Devign~\citep{zhou2019devigneffectivevulnerabilityidentification} established a framework that integrates abstract syntax, control flow, and data flow into a composite graph, utilizing gated graph propagation and a convolutional readout. ReVeal~\citep{chakraborty2020deeplearningbasedvulnerability} maintained the propagation mechanism but introduced a new readout, along with resampling and a representation-learning stage. Subsequent research varies either the graph structure or the aggregation method, drawing on the aggregation families catalogued in the graph-learning literature~\citep{zhou2021graphneuralnetworksreview}. For instance, LIVABLE~\citep{wen2023livable} implements differentiated propagation and a re-weighting loss for long-tailed type classification; Snopy~\citep{cao2024snopy} applies sample denoising and a causality-aware attention mechanism; DeepDFA~\citep{steenhoek2023dataflowanalysisinspireddeeplearning} abstracts the graph to dataflow facts and achieves strong results with substantially reduced training cost. Later work extends this in further directions: VulChecker~\citep{10.5555/3620237.3620604} localises defects within the graph rather than classifying whole functions, and counterfactual and edge-aware variants~\citep{10.1145/3650212.3652136,haque2025explainablevulnerabilitydetectioncc} alter what message passing attends to, while others vary the training setting rather than the model, as in federated training over code property graphs~\citep{ZHANG2024107371} or attention over data-dependence subgraphs~\citep{zhou2022grapheyenovelsolutiondetecting}. Across this work the encoding of the graph is treated as given: we are not aware of a study that varies edge direction or relation labelling and reports what either is worth.

The central premise is that program structure conveys information not present in flat token sequences. However, this assumption is generally asserted rather than empirically validated. Existing studies typically report end-to-end performance comparisons without isolating the marginal contribution of the graph within a controlled pipeline. Section~\ref{sec:graph-worth} presents such an empirical measurement.

\subsection{Token-based detection}
An alternative approach omits graph representations entirely. LineVul~\citep{fu2022linevul} fine-tunes CodeBERT on raw function text and reports significant improvements over graph-based predecessors, and subsequent re-evaluations find pretrained token models to be stronger baselines than the graph literature typically credits~\citep{li2025revisitingpretrainedlanguagemodels}. DeepDFA~\citep{steenhoek2023dataflowanalysisinspireddeeplearning} uses this method as the token baseline and finds the two approaches to be complementary. Pretrained code encoders such as GraphCodeBERT~\citep{guo2021graphcodebertpretrainingcoderepresentations}, which already incorporate data-flow structure in their objectives, are also commonly used as sources of node features for the aforementioned graph models. Consequently, the distinction between the two approaches is less clear-cut than it initially appears. When a graph model utilizes GraphCodeBERT embeddings, some of the structural information attributed to the graph is already provided by the encoder.


\subsection{What these models are learning}
Several lines of evidence suggest graph-based detectors rely on features other
than the program semantics they are assumed to capture. Interpretability studies
find that explanations produced for GNN detectors frequently do not align with
the statements a human would hold responsible~\citep{10.1145/3597926.3598145}.
Work on vulnerability-preserving transformation shows that predictions can be
flipped by edits that leave the defect
intact~\citep{10.1145/3652032.3657564}, and CausalVul~\citep{rahman2024causal}
removes variable names by construction because models are known to key on them.
Our measurements are consistent with this line of work: what the model extracts
from the graph separates classes in aggregate without localising the defect that
distinguishes a function from its own patch (Section~\ref{sec:mechanism}).

\subsection{Benchmarks and their defects}
Progress has been measured largely on a handful of corpora: Devign's FFmpeg+QEMU set~\citep{zhou2019devigneffectivevulnerabilityidentification}, Big-Vul~\citep{10148758,10.1145/3379597.3387501}, CVEfixes~\citep{10.1145/3475960.3475985,bhandari2021:cvefixes}, DiverseVul~\citep{chen2023diversevulnewvulnerablesource} and
MegaVul~\citep{Ni_2024}, with more recent efforts targeting collection quality directly~\citep{wang2023reefframeworkcollectingrealworld}. Each is mined from fix commits, and each inherits the labelling assumption that lines removed by a security patch were vulnerable.

The limitations of these benchmarks are extensively documented. ReVeal~\citep{chakraborty2020deeplearningbasedvulnerability} identified issues including duplication, unrealistic class balance, and models learning dataset artifacts rather than genuine vulnerability semantics. PrimeVul~\citep{ding2024primevul} addressed these concerns by deduplicating and re-labeling the data, followed by re-evaluating the same models, which led to a marked decline in reported performance. In paired evaluations that eliminate project-level cues by comparing a function with its patched version, model performance approaches random chance. Similarly, \citet{steenhoek2023empiricalstudydeeplearning} report instability across retraining sessions and datasets.


\subsection{Evaluation methodology beyond this task}
The concerns this paper measures are instances of patterns documented elsewhere,
and reading them together clarifies which of our findings are new.

That evaluation on data resembling the training set overstates deployment
performance was established in defect prediction two decades
ago~\citep{zimmermann2009cross}, and the same literature later showed that the
model validation technique alone shifts reported performance enough to change
which model appears
best~\citep{tantithamthavorn2017empirical}. In models of code more broadly,
near-duplicate functions are abundant in the repositories these corpora are
mined from~\citep{lopes2017dejavu} and measurably inflate results when they cross
the split~\citep{allamanis2019adverse}. Outside software engineering entirely,
leakage has been catalogued as a recurring cause of irreproducible findings
across machine-learning-based science~\citep{kapoor2023leakage}.

A parallel concern applies to the graph models themselves. Reported gains for
graph neural network architectures narrow substantially when model selection and
splitting are held fixed across the compared systems~\citep{errica2020fair}, and
guidance on assessing randomised algorithms in software engineering has long
recommended repeated runs with a test of significance rather than a single
measurement~\citep{arcuri2011practical}. Both bear directly on ablations whose
effect sizes are close to seed noise, which we show in
Section~\ref{sec:variance-result} is the regime this task occupies.

What this paper adds to that body of work is a decomposition rather than a
restatement: the leakage effect measured by type on one corpus with the model
held fixed, and the same apparatus turned on the graph encoding itself.

\subsection{Evaluation protocol}
The most closely related work comprises the critical literature discussed above, with generally convergent findings. The primary distinction lies in the specific aspects measured. ReVeal and PrimeVul demonstrate that performance declines under more realistic evaluation conditions; however, neither study decomposes this decline by leakage type while holding corpus, model, features, and random seed constant. CausalVul~\citep{rahman2024causal} addresses the issue differently by applying do-calculus to remove spurious features such as variable names, based on the observation that models frequently rely on such features.

Two gaps remain, both of which are addressed in this paper. First, no prior work quantifies the individual impact of each type of leakage (such as duplicates and patched twins, shared commits, and shared repositories), which is essential information for practitioners designing dataset splits. Second, and more importantly, no previous research isolates the contribution of the code property graph itself. Existing comparisons evaluate systems that differ simultaneously in representation, encoder, training regime, and hyperparameters. Controlling for all these factors and removing only the graph edges yields a different and less favorable assessment than system-level comparisons suggest.